\section{Data and Representations}\label{sec:data}

\subsection{Datasets and splits}
Two datasets are used. \textbf{GTZAN} \cite{tzanetakis2002gtzan} (Task 2): 10 genres, 999
usable 30-second clips (699/150/150 train/val/test) -- a fixed, study-specific, genre-stratified
split generated with a documented seed, not an official GTZAN benchmark split, used
throughout this study.
\textbf{MusicCaps} \cite{agostinelli2023musiclm} (Tasks 1, 3, 4): natural-language captions
paired with short tag phrases, used as a text-only corpus (Task 1) and, where audio is also
available, as caption-plus-audio-graph pairs
(Tasks 3, 4) -- these are two different subsets of MusicCaps, not one, since Task 1 needs no
audio and so can use every caption with usable text regardless of whether that clip's audio
was ever downloaded, while Tasks 3/4 are restricted to the smaller subset with a cached audio
graph. Task 1 uses 3,523/755/756 captions (756 test); the audio-paired subset used by Tasks
3 is 3,287/704/715.

\textbf{An important disclosure.} Tasks 1--3 share one implementation, though not one split
or segment length: Task 1 uses the caption-only 3,523/755/756 split above with no audio
segmentation; Task 2 uses GTZAN with 2-second segments; Task 3 uses MusicCaps with 1-second
segments and the audio-paired 3,287/704/715 split (50-tag vocabulary). Task 4 was built and
run separately, in a second, independent implementation of the same architecture family
(2-second MusicCaps segments; a 3,136/673/673 split assembled independently from this
implementation's own audio cache, not by augmenting the primary split above -- overlap
between the two subsets was not measured; a 30-tag vocabulary for its own zero-shot check).
Both implementations share the same node-feature definition, similarity-graph rule, and
encoder architectures (Section~\ref{sec:models}), and the cross-task argument
(Section~\ref{sec:cross-task}) does not depend on segment length or split size being
identical -- but the two are not byte-for-byte the same pipeline, stated plainly rather than
implying one implementation produced all four tasks' numbers.

\subsection{Audio segmentation and node features}
Task 1 uses no audio at all -- text only. Where audio is used (Tasks 2--4), each track is
split into non-overlapping segments: 2 seconds for GTZAN (Task 2) and for Task 4's MusicCaps
pipeline; 1 second for Task 3's MusicCaps pipeline (see the disclosure above). Each segment's
node feature vector is 12-dimensional chroma concatenated with 20-dimensional MFCC (32
dimensions total) in every audio pipeline, a purely local, per-segment summary with no
cross-segment information built in.

\subsection{Graph construction}
Edges connect (a) temporally adjacent segments, and (b) any pair of segments whose feature
cosine similarity exceeds a fixed threshold $\tau=0.8$ -- the same pairwise-segment-similarity
computation used in self-similarity-matrix analysis of audio structure \cite{foote2000novelty},
here thresholded into a graph's edge set rather than read off as a novelty curve. Writing
$X=\{x_1,\dots,x_n\}$ for a track's segment features, the similarity edge set is
\[
A_{\mathrm{sim}} = \{(i,j) : \cos(x_i,x_j) > \tau\},
\]
a deterministic function of $X$ alone -- a property that matters directly for interpreting
Section~\ref{sec:cross-task}. This scoping applies to similarity edges specifically, not the
deployed construction's temporal-adjacency edges, which encode segment order -- not
recoverable from an unordered feature set, outside this argument's premise. Measuring the
deployed construction directly confirms it is dense: mean edge density 95.4\% (MusicCaps) to
97.9\% (GTZAN), with 76.4\%-84.0\% of individual graphs \emph{fully} connected (a stricter,
related statistic; a companion audit measures a third dataset, MagnaTagATune, not used here,
at a comparably high 98.9\% density, not included in this range). The MusicCaps figure was
measured on Task 3's 715-example test split (1-second segments, roughly 10 nodes per graph);
not separately re-measured for Task 4's independent 2-second-segment pipeline, so we report
it as evidence about Task 3's graphs only. Randomly rewiring or resampling its edges
reproduces almost the same graph under a different name. A repaired, sparser construction
(mutual $k$-NN, $k\in\{4,8,16\}$, verified non-degenerate before use, though only three of
six dataset/$k$ combinations pass -- $k{=}16$ fails on both datasets and MusicCaps $k{=}8$
also fails) is used in a companion audit of this same system \cite{tismir_audit}; this paper
uses the original, denser construction throughout, since Task 2/3/4's checkpoints were
trained on it, noted as a limitation (Section~\ref{sec:discussion}).

\subsection{Text encoding}
Captions and tag phrases are tokenized with a standard BERT-family tokenizer and encoded
with a fine-tuned BERT/DistilBERT model (Section~\ref{sec:models}).

\subsection{Leakage and duplicate checks}
We ran three checks before treating any result as evidence.

\textbf{Caption-level label leakage (Task 1).} MusicCaps captions were checked for direct
lexical overlap with the target tag vocabulary: averaged across the 50 tags (label-macro, not
per-clip), 69.1\% of the time a tag's own word or a near-literal match appears in the caption
of a clip it labels. This is real and substantial, discussed directly in
Section~\ref{sec:cross-task} -- it explains part of why text alone is so strong in Task 1,
though a caption-sanitization control (target words removed, retrained from scratch) still
recovers 84.3\% of original performance, showing the effect is real but not total.

\textbf{GTZAN duplicate and split-leakage audit.} We hashed all 1,000 raw GTZAN audio files
(999 used across the splits) and found \textbf{14 exact, byte-identical duplicate pairs} --
consistent with GTZAN's documented quality issues \cite{sturm2013gtzan}. Cross-checking
against the split used throughout found \textbf{7 of the 14 pairs cross a split boundary}
(e.g.\ one copy in train, the other in test), meaning a small number of test tracks have a
byte-identical duplicate in the training set. One pair additionally carries two different
genre labels (\texttt{metal.00058.wav} and \texttt{rock.00016.wav} are byte-identical but
filed under different genres), a labeling inconsistency independent of splitting. In absolute
terms this affects a small fraction of the test set (\textbf{3 of 150 test tracks} have an
exact training-set duplicate). Unlikely to explain Task 2's largest gaps (CNN-vs-GraphSAGE
0.168, or the pooling-confounded edge-free comparison 0.045), but not ruled out for the
smaller, pooling-matched contrast (0.0072) -- see Section~\ref{sec:cross-task} for why. Real,
previously undocumented for this split, disclosed rather than corrected after the fact
(correcting would require retraining every Task 2 model on a new split, changing every
number here -- left as a direct next step, not silently patched).

\textbf{MusicCaps split-leakage audit.} All MusicCaps train/val/test splits were checked for
duplicate or overlapping YouTube IDs, both within a split and across splits. None were found
-- 3,287/704/715 unique IDs with zero overlap. This check is ID-based; it would not catch two
different clips of the same underlying recording, and audio-content-level near-duplicate
detection for MusicCaps was not run (Section~\ref{sec:discussion}).

\textbf{Artist-level leakage.} Neither dataset ships usable artist metadata (MusicCaps'
\texttt{author\_id} field identifies the caption's human annotator, not the track's
performing artist), so artist-level leakage could not be checked for either dataset. This is
a genuine, undisclosed-by-necessity limitation, not a check that was skipped.
